\section{Ambiente e tecnologias}
\label{sec:ambiente}

\paragraph{Distribuição Linux.} Existem distribuições próprias para
clusters, como a \textbf{\mbox{PelicanHPC}}~\cite{pelicanhpc}, que traz o MPI
pronto e dá boot nos workers pela rede. Porém, uma distribuição de cluster
não faz o paralelismo nem divide os dados, tarefa do programa do grupo: ela
apenas deixa pronta a configuração de que o MPI precisa (o mesmo MPI em
todos os nós, rede, SSH sem senha entre eles e a lista de máquinas). Como
essa configuração pode ser automatizada, o grupo optou pelo
\textbf{Ubuntu Server}, uma distribuição de uso geral, como costuma
ocorrer em clusters reais. Isso dá mais controle sobre cada peça e torna a
montagem do cluster parte do aprendizado.

\paragraph{Ambiente por etapa.} Docker e VMs são práticos, mas executam
tudo sobre o mesmo hardware: com quatro VMs em um notebook de quatro
núcleos, o ganho medido não representa um cluster real. Por isso, cada
ambiente foi atribuído à etapa em que é mais útil:

\begin{table}[H]
  \centering
  \small
  \begin{tabularx}{\textwidth}{>{\raggedright}p{2.9cm}>{\raggedright}p{4.0cm}X}
    \toprule
    \textbf{Etapa} & \textbf{Ambiente} & \textbf{Objetivo} \\
    \midrule
    Desenvolvimento & Máquina local ou Docker &
      Escrever e depurar o código \\
    Validação & VMs com Ubuntu Server, sem interface gráfica &
      Testar com rede e SSH reais e simular falhas desligando uma VM \\
    Medição & Notebooks físicos do grupo em rede &
      Medir o \emph{speedup} real, sem disputa de CPU \\
    \bottomrule
  \end{tabularx}
  \caption{Ambiente de cada etapa do projeto.}
\end{table}

\paragraph{Automação.} O \textbf{Vagrant} cria as VMs, e o
\textbf{Ansible} configura os nós (MPI, chaves SSH, lista de máquinas e
cópia do programa). O mesmo \emph{playbook} do Ansible configura as VMs e
os notebooks, bastando trocar a lista de máquinas, o que garante a mesma
configuração na validação e na medição. O grupo já usou o Ansible em
outras disciplinas.

\paragraph{Comunicação.} A tecnologia escolhida é o \textbf{MPI}
(\emph{Message Passing Interface}), padrão de troca de mensagens em
computação paralela, disponível em Python pela biblioteca
\texttt{mpi4py}. Ele oferece as operações \texttt{scatter} e
\texttt{gather}, que espalham os pedaços e recolhem os resultados, e
mensagens ponto a ponto para a fila dinâmica. Sockets, gRPC e filas de
mensagens (RabbitMQ) foram considerados, mas o MPI é o padrão da
computação em cluster e já cobre o que o modelo precisa.
